\chapter*{Introduction générale}
\addcontentsline{toc}{chapter}{Introduction générale}
\markboth{Introduction générale}{Introduction générale}

Les applications modernes reposent sur des serveurs, des conteneurs et des services qui
doivent fonctionner en continu. Lorsqu'un processeur sature, qu'un conteneur manque de
mémoire ou qu'un disque se remplit, l'équipe doit le savoir rapidement, comprendre la cause
et agir avant que les utilisateurs ne soient touchés. C'est le rôle de la
\textbf{supervision} (en anglais \emph{monitoring}) et de l'\textbf{alerte}.

Les outils de référence existent : Prometheus collecte les métriques, Alertmanager envoie
les alertes et Grafana les affiche. Ils sont puissants, mais ils forment trois produits
distincts à installer, configurer et relier entre eux. Pour une petite équipe, une startup ou
un projet étudiant, cette mise en place prend plusieurs jours. J'en ai fait l'expérience
lors de mon projet de fin d'année : la partie supervision a été la plus longue et la plus
fragile à configurer. De plus, ces outils déclenchent des alertes sur des seuils fixes, sans
expliquer pourquoi une alerte s'est déclenchée ni regrouper les alertes liées.

C'est dans ce contexte que s'inscrit mon stage au sein de \textbf{\entreprise}. Le projet,
nommé \textbf{SentinelOps}, a pour objectif de concevoir et réaliser une plateforme de
supervision et d'alerte \textbf{auto-hébergée}, déployable en \textbf{une seule commande},
qui collecte les métriques des serveurs et des conteneurs, détecte les incidents par des
règles et par une analyse statistique, explique chaque alerte avec des preuves mesurées et
notifie les équipes de manière fiable, le tout avec un niveau de sécurité adapté à un
environnement de production.

Ce rapport est organisé en cinq chapitres :
\begin{itemize}
  \item le \textbf{chapitre 1} présente l'organisme d'accueil, le contexte du projet,
        l'étude des solutions existantes et la méthodologie de travail ;
  \item le \textbf{chapitre 2} spécifie les besoins fonctionnels et non fonctionnels, ainsi
        que les acteurs et les cas d'utilisation ;
  \item le \textbf{chapitre 3} décrit la conception : architecture, modèle de données,
        automate d'états des alertes, détection d'anomalies et principales interactions ;
  \item le \textbf{chapitre 4} présente la réalisation : environnement, choix
        technologiques, implémentation des modules clés et interfaces ;
  \item le \textbf{chapitre 5} traite de la sécurité, de la conteneurisation, de la chaîne
        d'intégration continue, des tests, du déploiement et de l'évaluation.
\end{itemize}
Une conclusion générale dresse le bilan du travail et présente les perspectives.
